% ---------------------------------------------------------------------------
% Natural-Language Supervision of a 20-Robot Swarm
% arXiv preprint — DRAFT
%
% Build:  latexmk -pdf main.tex
%
% NOTE ON DOCUMENT CLASS
% ----------------------
% The predecessor report (paper/final.pdf) used IEEEtran. IEEEtran.cls is not
% installed in this container. To switch back once texlive-publishers is
% available, replace the \documentclass line with
%     \documentclass[conference]{IEEEtran}
% and delete the \authorblock helper definitions below.
% ---------------------------------------------------------------------------

\documentclass[10pt,twocolumn]{article}

\usepackage[margin=0.75in]{geometry}
\usepackage[T1]{fontenc}
\usepackage[utf8]{inputenc}
\usepackage{times}
\usepackage{booktabs}
\usepackage{amsmath,amssymb}
\usepackage{graphicx}
\usepackage{url}
\usepackage[hidelinks]{hyperref}
\usepackage{xcolor}
\usepackage{enumitem}
\usepackage{caption}
\usepackage{makecell}

\captionsetup{font=small,labelfont=bf}
\setlist{nosep,leftmargin=1.2em}

% Draft-note macro — remove before submission.
\newcommand{\todo}[1]{\textcolor{red}{\small[\textbf{TODO:} #1]}}
\newcommand{\pending}[1]{\textcolor{blue}{\small[\textbf{PENDING RESULTS:} #1]}}

\newcommand{\code}[1]{\texttt{\small #1}}

\title{\vspace{-2em}\bfseries Natural-Language Supervision of a 20-Robot Swarm:\\
An LLM Orchestration Layer over Classical MAPF\\ and Formation Control}

\author{
Evgenii Shlomov \quad Artyom Tuzov \quad Senya Belyh \quad Dmitry Ryabov\\
Dmitriy Vizitei \quad Anastasia Malakhova\\[0.3em]
\small Innopolis University, Innopolis, Russia\\
\small \texttt{\{e.shlomov, a.tuzov, s.belykh, d.ryabov,}\\
\small \texttt{d.vizitei, a.malakhova}@innopolis.university}
}
\date{}

\begin{document}
\maketitle

% ===========================================================================
\begin{abstract}
We present \code{iros\_llm\_swarm}, a complete ROS~2 system in which a large
language model supervises a 20-robot differential-drive swarm from free-form
natural language. Unlike prior LLM-robotics work that emits low-level code or
single-robot skill sequences, our orchestration layer emits a \emph{typed,
recursive plan tree} over five fleet-level primitives, and every leaf is
discharged by a classical algorithm --- Multi-Agent Path Finding (MAPF) with
two interchangeable backends (PBS and LNS2), a centralized leader--follower
formation manager, and a BehaviorTree.CPP mission front-end. The LLM therefore
never touches velocities, trajectories, or conflict resolution; it is confined
to a small, well-typed decision surface that the underlying stack can validate,
reject, and repair.

The contribution of this paper is the orchestration layer and its evaluation
methodology. We describe three control channels (reactive, proactive, and
operator chat), a grounding stack that gives the model bounded read-only access
to live fleet state through eleven semantic tools, schema-constrained decoding
against the plan grammar, and---centrally---a \emph{defence-in-depth} design in
which four deterministic guards and three LLM-in-the-loop repair loops sit
between a generated plan and the robots. We further contribute a reproducible
evaluation design: a 44-case hermetic planning benchmark spanning twelve
structural categories, and a proximity-driven task system that
turns end-to-end mission success into an automatically scored quantity, so that
planning accuracy and mission completion can be measured separately and
attributed to individual architectural components by ablation.

\pending{This draft fixes the system description and the experimental design.
E1 (hermetic planning accuracy) is populated across five local models;
E2--E6 remain unpopulated pending further runs.}
\end{abstract}

% ===========================================================================
\section{Introduction}

Coordinating a fleet of mobile robots that share one workspace is a long-standing
problem. The classical decomposition separates fleet-level decisions---which
robot goes where, in what order---from robot-level execution---how one robot
follows a path or holds an offset. Large language models make this decomposition
newly attractive~\cite{saycan,codeaspolicies,progprompt,innermonologue}: an LLM
can reason over compact semantic primitives instead of producing control
commands, while the classical stack retains responsibility for safety, conflict
resolution, and real-time tracking.

The difficulty is that this promise is easy to state and hard to make hold. A
language model asked to command twenty robots will confidently emit coordinates
inside walls, address robots that are currently locked into a formation, place
four robots on the same point, invent room names, and return JSON that does not
parse. Any one of these silently produces a fleet-wide failure. The engineering
question this paper addresses is therefore not \emph{can an LLM plan for a
swarm} but \emph{what has to sit between the LLM and the robots so that it
can}, and how much each of those pieces is actually worth.

\subsection{Scope and contributions}

We study a fleet of 20 differential-drive robots in a 2D environment simulated
with Stage, each running its own Nav2 stack. The system spans three layers: a
natural-language orchestration layer (this paper's focus), a classical
coordination substrate (MAPF, formations, behavior trees), and per-robot
navigation. Our contributions:

\begin{enumerate}[label=(\arabic*)]
\item \textbf{A typed plan language for fleet control.} Five leaf primitives
(\code{mapf}, \code{formation}, \code{disband}, \code{idle}) and two containers
(\code{sequence}, \code{parallel}) form a recursive tree that is expressive
enough for multi-group, multi-phase missions yet small enough to be validated
exhaustively and to be expressed as a JSON schema for constrained decoding
(Sec.~\ref{sec:planlang}).

\item \textbf{Three orchestration channels} over one substrate: a reactive
channel invoked by behavior-tree nodes on WARN/ERROR, a proactive channel that
observes fleet state and intervenes on its own initiative, and an operator chat
channel that streams a reply and an executable plan (Sec.~\ref{sec:channels}).

\item \textbf{A grounding stack} that replaces coordinate hallucination with
lookup: a live pose cache, a room-occupancy-aware group placement solver, and
eleven read-only semantic tools brokered behind an allowlist, with an optional
MCP provider for raw ROS introspection (Sec.~\ref{sec:grounding}).

\item \textbf{Defence in depth between plan and robots}: schema-constrained
decoding, permissive type coercion, an active-formation guard, an
occupancy-aware goal rewriter, and automatic formation staging---four of which
are deterministic and none of which require another model call
(Sec.~\ref{sec:guards}).

\item \textbf{A bounded closed loop}: deterministic post-execution
verification feeding three nested LLM repair loops (remediation, verification
repair, mission supervision), each with an explicit step, time, and
no-progress budget (Sec.~\ref{sec:closedloop}).

\item \textbf{A 20-robot execution substrate} that makes the above runnable and
measurable: a Stage 2D environment reached after a documented pivot from Gazebo
Harmonic, five scenarios with symbolic map descriptions and a dynamic obstacle
system, per-robot Nav2 stacks generated from a single parameter template with
staggered bring-up, and a static-localization strategy that replaces AMCL at
this fleet density (Sec.~\ref{sec:substrate}).

\item \textbf{A reproducible evaluation design} (Sec.~\ref{sec:expdesign}) that
separates planning competence from execution competence, scores end-to-end
missions automatically via a proximity-driven task system, and attributes
outcomes to individual components through configuration-level ablation.
\end{enumerate}

The system was previously documented as two partial technical reports written
against a three-level split of the work: one covering the classical
coordination substrate~\cite{ourlevel2}, the other the simulation environment,
per-robot navigation, and an early design of the orchestration
layer~\cite{secondteam}. Each necessarily described its layer against
\emph{assumed} interfaces at the boundary. This paper is the unified account of
the system as built: it supersedes the evaluation sections of both reports,
corrects the orchestration-layer description to match the shipped
implementation, and---because planning quality and mission success can only be
compared on the whole stack---evaluates the three layers together rather than
separately.

% ===========================================================================
\section{Related Work}

\paragraph{Multi-Agent Path Finding.}
MAPF has a rich algorithmic literature~\cite{stern}. Conflict-Based
Search~\cite{cbs} is the canonical optimal solver. Priority-Based Search
(PBS)~\cite{pbs} relaxes optimality for scalability: a strict priority ordering
replaces time-indexed constraints, and each agent plans deterministically
against higher-priority paths. Large Neighborhood Search for MAPF
(LNS2)~\cite{lns2} is anytime and suboptimal, starting from a possibly
colliding plan and iteratively destroying and repairing agent subsets with
soft-constraint A*. We adopt both behind one action interface, which lets us
treat the backend as an experimental factor rather than a design commitment.

\paragraph{Formation control.}
Formation control is typically posed as leader--follower, virtual structure, or
behavior-based~\cite{formsurvey1,formsurvey2}. We use leader--follower because
it composes cleanly with a per-robot motion controller and matches the kind of
semantic offset (``a wedge'', ``abreast'') a language model naturally produces.

\paragraph{Navigation and behavior trees.}
Nav2~\cite{nav2} is the de facto ROS~2 navigation stack, with layered
costmaps~\cite{layered} the standard obstacle representation. Behavior
trees~\cite{btbook} via BehaviorTree.CPP~\cite{btcpp} provide a modular control
architecture that is particularly suitable when an external reasoning component
supplies tree parameters.

\paragraph{LLMs for robotics.}
SayCan~\cite{saycan} grounds language in affordances; Code as
Policies~\cite{codeaspolicies} and ProgPrompt~\cite{progprompt} emit executable
program text; Inner Monologue~\cite{innermonologue} closes the loop with
textual feedback. Our design differs in two respects that matter at fleet
scale. First, the model emits a \emph{typed tree over fleet primitives}, not
code: the output space is small enough to constrain by schema and validate
exhaustively, which removes the arbitrary-code-execution surface and the
malformed-output failure mode. Second, feedback is not free text but a
\emph{deterministic verification record} computed from formation status and
pose state (Sec.~\ref{sec:closedloop}), so the repair loop is driven by
measured ground truth rather than by the model's own narration.

\paragraph{ROS~2 and DDS at scale.}
DDS discovery cost dominates at scale~\cite{dds1,dds2}. We use CycloneDDS with
a project-specific profile; without it the 20-robot system does not reliably
complete discovery within the launch budget.

% ===========================================================================
\section{The Classical Substrate}
\label{sec:substrate}

This section summarizes the layer the orchestrator commands; see~\cite{ourlevel2}
for full detail.

\paragraph{Shared fleet action.}
Both MAPF backends own the same long-lived \code{/swarm/set\_goals} action,
which encapsulates the entire
$\textit{validating} \rightarrow \textit{planning} \rightarrow
\textit{executing} \rightarrow \textit{replanning} \rightarrow
\textit{succeeded}$
lifecycle behind one goal token. Schedule monitoring and replanning are
internal; clients see one atomic operation. Cancellation stops every robot. The
action result carries a full metrics payload---planning time, node expansions,
per-agent path lengths, replan count, total execution time---which we exploit
directly for benchmarking (Sec.~\ref{sec:e4}).

\paragraph{PBS backend.}
An 8-connected A* on a grid downsampled to 0.2\,m, with a backward-Dijkstra
heuristic per unique goal, two-zone gradient inflation (hard footprint plus a
soft penalty zone), and capsule-vs-capsule conflict detection that subsumes the
usual vertex/edge split. The high-level search follows the priority-based
formulation of~\cite{pbs}.

\paragraph{LNS2 backend.}
An initial soft-constraint A* solution refined by an adaptive
large-neighborhood search with four destroy operators (collision-based,
random-walk, random, bottleneck) whose weights adapt online. It is anytime,
tolerant of transiently infeasible starts, and warm-starts on replan.

\paragraph{Per-robot followers and formations.}
Each robot runs a dual-mode follower: in AUTONOMOUS mode it dispatches its
assigned plan to Nav2 \code{follow\_path}; in FORMATION mode it disengages Nav2
and runs a PD controller tracking
$p_{\mathrm{ref}} = p_{\mathrm{leader}} + R(\theta_{\mathrm{leader}})\,o$
for a body-frame offset $o$. Mode switching is driven by a single latched
message, without node restarts. A centralized manager owns formation
definitions and republishes the full registry on a latched topic; a monitor
computes per-follower error
$\varepsilon_i = \lVert p_i - (p_{\mathrm{leader}} + R(\theta_{\mathrm{leader}}) o_i) \rVert$
and classifies each formation as INACTIVE / FORMING / STABLE / DEGRADED /
BROKEN at 10\,Hz. These $\varepsilon_i$ are published, which makes formation
tracking accuracy directly measurable (Sec.~\ref{sec:e5}).

\paragraph{Behavior-tree front-end.}
BehaviorTree.CPP v3 nodes wrap every primitive: \code{MapfPlan},
\code{SetFormation}, \code{DisableFormation}, and a \code{CheckMode} guard.
Mission completion is observed on the \code{/bt/state} blackboard mirror, whose
\code{active\_action} field cycles \code{none} $\rightarrow$ node $\rightarrow$
\code{none}; the orchestrator keys leaf completion on this transition rather
than on the command acknowledgement, which is only an ack.

\paragraph{Simulation environment.}
The fleet runs in Stage 2D~\cite{stage} via \code{stage\_ros2}, with a unified
TF tree and enforced namespace prefixes. An earlier Gazebo
Harmonic~\cite{gazebo} configuration was implemented end-to-end and then
retired: per-robot physics, collision broadphase, and sensor rendering pushed
the real-time factor below usable thresholds at 20 robots before any navigation
stack was attached~\cite{secondteam}. The Gazebo configuration is retained in
tree for smaller fleets and future hardware-realistic preview.

Five scenarios share one robot model and spawn protocol: \textsf{cave}
(75$\times$75\,m, narrow single-file passages), \textsf{large\_cave}
(150$\times$150\,m, same topology at double scale), \textsf{warehouse\_2} and
\textsf{warehouse\_4} (30$\times$30\,m, two- and four-group symmetric layouts),
and \textsf{amongus} (66.6$\times$37.75\,m), which maps the spaceship layout of
\emph{Among Us: The Skeld} into 14 named rooms joined by narrow corridor
choke-points. We use \textsf{amongus} for all orchestration experiments because
its rich naming exercises symbolic reference and its choke-points stress
conflict resolution.

Each scenario ships a machine-readable map description consumed by the
orchestrator: named locations with coordinates, aliases in English and Russian,
colour-coded robot groups with home and spawn positions, formation-capable
zones with free-space radii, and map-specific planning heuristics (bottleneck
locations, known-unreachable goal patterns, abort-versus-replan guidance). A
dynamic obstacle manager additionally supports runtime-togglable doors and
placeable obstacles; \textsf{amongus} declares six corridor doors matching the
map geometry.

\paragraph{Scale engineering.}
Three decisions are load-bearing at 20 robots and are reported because they
constrain the experiments. CycloneDDS with a custom profile is mandatory above
roughly ten participants. AMCL is disabled---each robot's scan contains every
other robot's body, and particle weights collapse---so localization is a static
per-robot \code{map} $\rightarrow$ \code{odom} identity transform over simulator
odometry. Bring-up is staggered (simulator 0\,s, Nav2 1\,s, planner 10\,s,
followers 12\,s, plus 0.3\,s per Nav2 instance).

% ===========================================================================
\section{Orchestration Layer}

\subsection{Three channels}
\label{sec:channels}

The orchestrator exposes three independent paths from language to fleet action,
which share prompt infrastructure and the plan validator but differ in trigger,
latency budget, and authority.

\textbf{Channel 1 (reactive).} A \code{/llm/decision} action server. Behavior-tree
nodes call it when a primitive reports WARN or ERROR, supplying the event plus a
tail of the accumulated log. The model returns exactly one of
\code{wait} / \code{abort} / \code{replan}. The prompt biases hard toward
\code{wait} and enumerates planner-internal warnings that indicate
\emph{self-recovery in progress} (``replan did not reach collision-free; keeping
old paths'', ``warm seed rejected'') as explicit \code{wait} cases---because
the expensive error here is a spurious \code{replan} that cancels a plan the
substrate was about to fix on its own. Always active.

\textbf{Channel 2 (proactive).} A passive observer subscribing to
\code{/bt/state}. It decides on its own initiative whether to intervene and
publishes an \code{LlmCommand} that a behavior-tree receiver applies to the
blackboard. It is cooldown-gated and checks an \code{llm\_thinking} flag to
avoid colliding with channel~1. Off by default.

\textbf{Channel 3 (operator chat).} The main path, and the focus of this paper.
An operator types free-form text; the server streams a natural-language reply
token-by-token, parses an executable plan from the same response, and dispatches
it. Everything in Sections~\ref{sec:planlang}--\ref{sec:closedloop} describes
channel~3.

\subsection{Plan language}
\label{sec:planlang}

Channel~3 returns a single JSON object \code{\{reasoning?, reply, plan\}}, where
\code{plan} is a recursive tree over the node types in
Table~\ref{tab:planlang}.

\begin{table}[t]
\centering\small
\caption{Plan tree node types.}
\label{tab:planlang}
\begin{tabular}{@{}llp{4.1cm}@{}}
\toprule
\textbf{Node} & \textbf{Kind} & \textbf{Payload / semantics} \\
\midrule
\code{mapf}      & leaf & \code{robot\_ids}, \code{goals}, \code{spread}; one goal per robot, or one centre that the server grids \\
\code{formation} & leaf & \code{formation\_id}, \code{leader\_ns}, \code{follower\_ns}, \code{offsets\_x/y} \\
\code{disband}   & leaf & \code{formation\_id} \\
\code{idle}      & leaf & \code{reason}; also the carrier for escalation (see below) \\
\code{sequence}  & container & children in order; abort on first failure \\
\code{parallel}  & container & see flattening rules below \\
\bottomrule
\end{tabular}
\end{table}

Three design decisions in this grammar are worth stating explicitly.

\emph{Parallelism is a claim about intent, not about threads.} A
\code{parallel} node is flattened before execution: multiple \code{mapf} leaves
\emph{merge into one} \code{mapf} goal, because the MAPF planner already
resolves inter-robot conflict internally and issuing two concurrent fleet goals
would defeat exactly the coordination we want. Mixed \code{mapf} and
\code{formation} children degrade to sequential execution, since a formation
requires its own follower mode. An \code{idle} in any branch short-circuits the
whole node---stop takes priority. This means the operator's ``do A and B at the
same time'' is honoured as \emph{one jointly planned motion}, which is stronger
than concurrent dispatch, not weaker.

\emph{\code{idle} doubles as the escalation channel.} Rather than adding node
types for ``I need clarification'' and ``here is an answer, no action needed'',
the model emits \code{idle} with a prefixed reason: \code{needs\_help:},
\code{clarify:}, or \code{reply\_only:}. This keeps the grammar at six node
types---small enough to constrain by schema---while still distinguishing
refusal-to-act from action. It also makes escalation a first-class, testable
behaviour: a benchmark case can assert that a question with a missing robot
pose produces \code{needs\_help:} and not a confident plan over stale data.

\emph{The schema is enforced twice.} \code{plan\_schema.py} expresses the
grammar as a JSON schema with a recursive \code{\$defs/node} reference, passed
to the backend for constrained decoding (Ollama \code{format}, OpenAI
\code{response\_format}). But JSON Schema cannot express the cross-field
constraint $|\code{goals}| = |\code{robot\_ids}|$, so a hand-written validator
re-checks every node after parsing. Constrained decoding removes the
\emph{malformed} failure mode; the validator removes the \emph{incoherent} one.

\subsection{Grounding}
\label{sec:grounding}

A model asked for ``cyan to the cafeteria'' must resolve a colour group to
robot ids, a room name to coordinates, and coordinates to free floor space. We
push all three out of the model's head and into lookup.

\emph{Static context.} The prompt carries a per-map description: named
locations with coordinates, colour groups with member ids and home positions,
formation-capable zones with radii, and door geometry. Maps are declared in
\code{common\_scenarios.yaml}, so adding an environment does not touch the
orchestrator.

\emph{Live context.} A pose cache subscribes to \code{/robot\_*/odom} and
supplies fresh positions with an explicit staleness bound (2\,s). This is what
lets the model compute formation goals as $\textit{leader}_{xy} + \textit{offsets}$
rather than inventing absolute coordinates---and, equally important, lets it
\emph{notice a missing robot} and escalate instead of planning over a gap.
Formation status, behavior-tree state, recent events, and open task state are
folded into the same snapshot under a character budget.

\emph{Semantic tools.} Eleven read-only tools sit above raw introspection:
system state, group state, location resolution, route context with bottleneck
and door penalties, goal and formation feasibility checks, occupancy-aware free
group-goal search inside a room, post-execution verification, the allowed
action schema, and recovery options. They are dispatched by a broker that
enforces an allowlist, a per-round tool cap, a timeout, and a result-size cap.
An optional MCP provider (\code{ros-mcp} over stdio) adds raw ROS introspection
under a separate allowlist restricted to read verbs. The model never sees MCP
tools directly---only a bounded snapshot or a brokered semantic call---so the
write surface is structurally absent rather than merely discouraged.

Notably, tool calling and schema-constrained decoding are \emph{mutually
exclusive} in current backends: constrained decoding suppresses tool-call
emission. The system therefore has two distinct planning modes, which we treat
as a primary experimental factor (Sec.~\ref{sec:e2}).

\subsection{Defence in depth}
\label{sec:guards}

Table~\ref{tab:guards} lists what stands between a generated plan and the
robots. The design principle is that \emph{a deterministic guard is always
preferred to another model call}: guards are cheap, reproducible, and cannot
themselves hallucinate.

\begin{table}[t]
\centering\small
\caption{Layers between a generated plan and the fleet. D = deterministic,
L = requires an LLM call.}
\label{tab:guards}
\begin{tabular}{@{}llp{3.5cm}@{}}
\toprule
\textbf{Layer} & & \textbf{Failure mode addressed} \\
\midrule
Constrained decoding      & D & unparseable / off-grammar JSON \\
Schema validator          & D & structurally valid but incoherent nodes \\
Type coercion             & D & \code{"robot\_10"} where an \code{int} is required \\
Active-formation guard    & D & direct MAPF control of a robot locked in a live formation \\
Occupancy rewrite         & D & room-centre goals that land in walls or clump \\
Formation prestaging      & D & bare \code{formation} leaf with followers out of position \\
\midrule
Remediation loop          & L & execution failed; replan with failure briefing \\
Verification repair       & L & execution reported success but the outcome does not hold \\
Mission supervision       & L & multi-step mission needs continuation \\
\bottomrule
\end{tabular}
\end{table}

Two of these deserve elaboration because they encode domain knowledge the model
does not reliably have.

\emph{Type coercion} is deliberately permissive. Small local models routinely
emit \code{"robot\_10"} or \code{10.0} where the grammar wants \code{10}. A
strict parser would reject an otherwise perfect plan over a namespace string.
The coercer accepts \code{10}, \code{"10"}, \code{"robot\_10"},
\code{"robot10"}, and integral floats, while explicitly rejecting booleans
(which are \code{int} subclasses in Python and would otherwise silently become
robot~0 or~1). Strictness is spent where it buys safety, not where it only
punishes formatting.

\emph{The active-formation guard} exists because ``send the magenta line to
storage'' is ambiguous in a way that matters. A formation is commanded by moving
its \emph{leader}; commanding all four members individually tears it apart. The
guard reads live formation status, and for any formation in FORMING, STABLE, or
DEGRADED it strips follower ids out of MAPF leaves, leaving the leader. It
returns a structured failure record shaped exactly like a PlanExecutor failure,
so the existing repair loop consumes it with no special handling.

\subsection{Closed loop}
\label{sec:closedloop}

Execution success is not the same as mission success, and this distinction
drives the whole loop design. A plan can execute cleanly---every action server
returns success---while the requested outcome does not hold: robots arrive but
the formation never stabilizes, or they arrive at a room's centre and pile up.
The system therefore verifies deterministically after execution, reading
formation status and pose state to decide whether what was asked for is
actually true, and emits an honest \emph{partial} result when state is missing
rather than guessing.

Three nested loops consume that verdict, each explicitly budgeted:

\begin{itemize}
\item \textbf{Remediation} (default 2 attempts) fires when execution
\emph{fails}. It re-prompts with the failed leaf, the error, the history of
prior attempts, and a \emph{targeted refresh} of runtime context---not the full
snapshot, to keep the retry prompt inside the context budget.
\item \textbf{Verification repair} (default 2 attempts) fires when execution
\emph{succeeds} but verification says the outcome does not hold, and only when
the verifier itself marks the discrepancy repairable.
\item \textbf{Mission supervision} (default 6 steps, 180\,s) wraps both in a
bounded execute--verify--continue loop for multi-phase missions, with a
no-progress limit of 2 and a fingerprint set of already-failed plans so the
model cannot loop on a plan that has already been tried.
\end{itemize}

Every loop terminates in a help escalation to the operator rather than in
silent failure or unbounded retry. We regard these budgets as tuning
parameters, not constants, and measure their effect directly
(Sec.~\ref{sec:e3}).

\subsection{Task system as mission ground truth}
\label{sec:tasks}

A separate node maintains a set of spatial tasks and polls TF at 5\,Hz. A
\emph{point} task transitions PENDING $\rightarrow$ DONE when any robot enters
its radius. A \emph{carry} task transitions PENDING $\rightarrow$ CARRYING on
pickup and $\rightarrow$ DONE when the carrying robot reaches the dropoff.
Tasks are declared per scenario in YAML and their state is published.

This component was built for mission planning, but it doubles as the
measurement instrument that makes end-to-end evaluation tractable: it converts
``did the swarm accomplish what the operator asked'' into an automatically
scored, timestamped, operator-independent quantity. Every end-to-end experiment
below is scored from this signal.

% ===========================================================================
\section{Experimental Design}
\label{sec:expdesign}

This section is the methodological core of the paper. We state it in full
before reporting numbers so that the measurement protocol is fixed independently
of its outcome.

\subsection{Research questions}

\begin{description}[leftmargin=2.6em,style=nextline]
\item[RQ1] How accurately does a locally-hosted open-weight LLM translate
free-form fleet commands into valid plans, and how does accuracy vary across
linguistic and structural categories?
\item[RQ2] What does each grounding and constraint mechanism contribute?
Specifically: schema-constrained decoding vs.\ tool calling, live pose context,
and context budget.
\item[RQ3] Does planning accuracy predict mission success, or does the gap
between them---the part closed by guards and repair loops---dominate?
\item[RQ4] What does each closed-loop layer contribute to end-to-end mission
completion, and at what cost in time and model calls?
\item[RQ5] How do the two MAPF backends compare under identical fleet-level
requests, and does backend choice interact with orchestration quality?
\end{description}

RQ3 is the question we consider most important and the one the literature most
often skips: a benchmark measured on plan text alone can look excellent while
the fleet still fails, and the converse---weak plans rescued by strong
guards---is equally possible. Measuring both on the same system is the point.

\subsection{Principle: reuse existing instrumentation}

A deliberate constraint on this design is that it introduces almost no new
measurement code. Every metric below is read from a signal the running system
already produces (Table~\ref{tab:metrics}). This matters for validity: an
evaluation harness that reimplements the notion of success is measuring the
harness, and instrumentation added for a paper tends not to survive the paper.

\begin{table}[t]
\centering\small
\caption{Metrics and their existing sources.}
\label{tab:metrics}
\begin{tabular}{@{}p{3.3cm}p{4.4cm}@{}}
\toprule
\textbf{Quantity} & \textbf{Source (already implemented)} \\
\midrule
Plan validity, category accuracy, latency & 44-case benchmark harness \\
MAPF planning time, expansions, path lengths, replans & \code{SetGoals} action \emph{result} payload \\
Arrival / deviation / stall counts & \code{SetGoals} action \emph{feedback} \\
Formation error $\varepsilon_i$, state machine & \code{FormationStatus} @ 10\,Hz \\
Mission completion, time-to-complete & \code{TaskState} from the task manager \\
Repair / remediation / supervision counts & channel-3 event stream \\
\bottomrule
\end{tabular}
\end{table}

\subsection{E1: Planning accuracy (hermetic, offline)}
\label{sec:e1}

The 44-case benchmark calls the model directly---no simulator, no ROS---with
\emph{fixed synthetic runtime contexts} (fleet at spawn; a stable magenta line;
a stable cyan wedge; robot~12's pose missing; robot~4 at the cafeteria). Each
case pairs a prompt with a programmatic validator asserting structural
properties of the returned tree: node type, exact robot-id sets, goal proximity
to a named location within a tolerance, goal distinctness, formation leader
identity, staging-before-formation ordering, and escalation-reason prefixes.

Cases span twelve categories chosen to separate distinct competences:
basic single-group \code{mapf}; explicit per-robot coordinates; goal spreading;
\code{parallel} over multiple groups; multi-phase \code{sequence}; formation
creation (which requires \emph{staging followers first and excluding the
leader}); formation movement (which requires commanding \emph{only} the
leader); \code{formation\_complex} (creation followed by a leader-only move);
disband; \code{idle}; escalation under missing or ambiguous state; and a set of
edge cases (repeated/ambiguous commands, partial formations, adjacency
requests). All prompts here are English. An earlier revision of this
benchmark additionally scored reply-language fidelity on a parallel set of
Russian-language prompts; we dropped that set because it tests a UI-level
concern (does the chat reply mirror the operator's language) orthogonal to
planning accuracy, and folding it into a single pass rate obscured both
signals.

\emph{Protocol.} Each case is run $N{=}5$ times per model at temperature 0.1.
We report pass@1 (mean over repeats) and pass@5, per-category breakdown,
end-to-end latency, and the rate of each distinct failure mode (unparseable,
schema-invalid, wrong ids, wrong coordinates, wrong structure). All runs
reported here use \code{--llm-max-tokens 2048} (raised from an initial
768, which was tight enough to make truncated JSON a confound for more
verbose models independent of their planning ability) and
\code{--llm-num-ctx 32768}.

\emph{Models.} The deployment target is a single workstation with 16\,GB VRAM,
which fixes the candidate set. The prompt is large---we measured
$\approx$9{,}700 tokens for the system prompt plus seven few-shot exchanges,
before runtime context---so KV-cache cost is a real constraint, not a rounding
error. We sweep five locally-hosted instruction models spanning that budget:
a small model (\textbf{qwen3.5:4b}, 3.4\,GB); a mid-size model from a
different family (\textbf{gemma3:12b}, 8.1\,GB); the intended 4-bit 14B
operating point (\textbf{qwen2.5:14b}, 9.0\,GB); an 8-bit $\approx$9B model
(\textbf{qwen3.5-9b}, 10\,GB); and the largest model that still fits under the
VRAM ceiling (\textbf{mistral-small3.2:24b}, 15\,GB, 4-bit, partial CPU
offload expected). Whether the last of these remains practical under offload
is itself part of what the latency column reports.

\emph{Threat.} E1 measures planning in isolation against synthetic context. It
deliberately cannot detect a plan that is structurally perfect and physically
unexecutable. That is E3's job, and the comparison between them answers RQ3.

\subsection{E2: Grounding and constraint ablation}
\label{sec:e2}

Same harness, same cases, toggling one mechanism at a time:
schema-constrained decoding on/off; live pose context on/off; goal auto-spread
on/off; and context budget swept across its range. All of these are existing
runtime parameters, so the ablation requires no code change.

We add one factor that is not currently a parameter: \emph{symbolic versus raw
grounding}. During development we observed that prompting with raw coordinates
produced a higher rate of out-of-bounds and wall-interior goals than symbolic
room references, and that supplying named locations, formation zones, and
map-specific heuristics substantially reduced goal-validation failures
downstream~\cite{secondteam}. That observation was uncontrolled and
retrospective. Because the map description is a separable prompt block, it can
be ablated cleanly---full description, locations without heuristics, and
coordinates only---turning an anecdote into a measured effect size. Given that
the map description is a large share of a $\approx$9{,}700-token prompt, this
also quantifies what that context is buying.

One asymmetry must be stated plainly. The offline harness invokes the model
directly and therefore exercises only the plain and schema-constrained paths;
the \emph{tool-calling} path lives in the chat server. Comparing constrained
decoding against tool calling---the mutual exclusivity noted in
Sec.~\ref{sec:grounding}, and arguably the single most interesting ablation
here---therefore requires either extending the harness to route through the
chat server, or measuring that factor only in E3. We take the first option, as
the second confounds planning mode with execution noise.

\subsection{E3: End-to-end missions}
\label{sec:e3}

The primary experiment. A 20-robot stack runs in the 14-location environment
with nine declared tasks (six point, three carry). An operator command is
issued through channel~3; the system plans, executes, verifies, and repairs;
the task manager scores completion independently.

\emph{Factors.} Remediation on/off; verification repair on/off; mission
supervision on/off; planning mode (constrained vs.\ tool-calling); model size.

\emph{Metrics.} Task completion rate; time to completion; number of LLM calls
per mission (the cost of each safety net); how often each guard fired; how
often each repair loop fired and whether it succeeded; escalation rate; and
total fleet distance travelled.

\emph{Protocol.} $N$ repeats per configuration with a fixed initial fleet
state, resetting task state between runs via the existing reset service.
Because a 20-robot stack is stochastic---DDS timing, Nav2 controller
noise---single runs are not informative, and we report distributions, not
point estimates. This directly addresses our earlier Level-2 report's
acknowledged limitation that its integration scenario was run once.

\emph{No-LLM baseline.} The workspace includes a scripted behavior-tree runner
that drives the identical six-step mission---fleet convergence, two formation
activations, a full diagonal cross-swap, return home, idle---against the same
live 20-robot stack with no LLM in the loop~\cite{secondteam}. This is
the natural ceiling for E3: it isolates what the substrate achieves when plans
are known-correct by construction. The difference between scripted completion
and LLM-driven completion is precisely the orchestration cost, and reporting
E3 without it would leave the reader unable to tell a weak orchestrator from a
hard environment.

\emph{Guard attribution.} Because guards are deterministic, their contribution
is measurable without disabling them: we log every firing and, for each,
whether the plan would have been physically invalid without it. Disabling a
safety guard on a live 20-robot fleet to measure it is neither necessary nor
wise.

\subsection{E3b: Channel-1 decision accuracy}
\label{sec:e3b}

Channels~1 and~2 are the least evaluated part of the system in either report,
and unlike channel~3 they already write structured per-call JSONL logs
(prompt, raw output, parsed verdict). This makes a replay evaluation nearly
free: accumulate incident records across the E3 and E4 campaigns, have two
annotators independently label the correct verdict
(\code{wait}/\code{abort}/\code{replan}) from the event and log tail, and score
the model against that label set with inter-annotator agreement reported.

The interesting quantity here is not overall accuracy but the
\emph{asymmetric} error rate. The prompt deliberately biases toward
\code{wait}, because a spurious \code{replan} cancels a plan the substrate was
about to repair on its own, whereas a spurious \code{wait} merely costs one
supervision cycle. We therefore report the confusion matrix rather than a
scalar, and separately the false-\code{replan} rate on the planner-internal
warnings the prompt enumerates as self-recovery cases.

\subsection{E3c: Recovery under environmental change}
\label{sec:e3c}

The repair loops of Sec.~\ref{sec:closedloop} exist for the case where the
world changes after a plan is made, yet every experiment above holds the
environment fixed. The dynamic obstacle manager closes this gap directly:
closing one of the six \textsf{amongus} corridor doors mid-mission invalidates
a committed route and forces the stack to notice, replan, and---if the goal
becomes unreachable---escalate.

We script door closures at a fixed mission phase and measure: whether the
MAPF layer replans without orchestrator involvement (the desired outcome, no
LLM call); whether the closure instead surfaces as a channel-1 WARN and which
verdict follows; time from closure to recovered motion; and mission completion
compared to the undisturbed run. A door closure that makes a task
\emph{unreachable} is the sharpest available test of escalation: the correct
behaviour is \code{needs\_help:}, and a confident replan onto an impossible
goal is a clear failure.

\emph{Caveat.} This experiment requires integration work the others do not.
The obstacle manager and the zone map server both publish the map with
transient-local durability and cannot currently coexist; the obstacle manager
is not wired into the bring-up launches. Composing it means repointing the
static map server to the raw map topic first. We flag this as a prerequisite
rather than assuming it away.

\subsection{E4: MAPF backend comparison}
\label{sec:e4}

PBS vs.\ LNS2 under identical fleet requests, using the existing goal-dispatch
test client: convergence to a single point, random goals within a radius, and
group-home dispersal, across 5/10/15/20 robots and three environments of
increasing clutter. All metrics come from the action result payload. This
closes our earlier Level-2 report's ``no statistics over repeated runs'' gap and
tests the interaction half of RQ5.

\subsection{E5: Formation tracking accuracy}
\label{sec:e5}

Recording the formation status topic while a leader traverses a fixed route
yields $\varepsilon_i(t)$ directly. We characterize steady-state and peak error
by shape (line, wedge, triangle, abreast) and leader speed, plus time-to-STABLE
after activation and DEGRADED dwell fraction. This is a bag recording and an
offline parse; our earlier Level-2 report listed it as an open gap purely
because nobody had recorded it.

\subsection{E6: Costmap layer ablation (secondary)}
\label{sec:e6}

Comparing the custom self-clearing costmap layer against the stock obstacle
layer under 20-robot scan density, measuring recovery-behaviour invocations and
mission success. Included for completeness; lowest priority.

\subsection{Instrumentation to be added}

Three small additions are required, and we note them explicitly because they
are the only new code the evaluation needs:

\begin{enumerate}[label=(\roman*)]
\item A per-mission JSONL logger on channel~3. Channels~1 and~2 already write
structured logs through a shared writer; channel~3 does not, and E3 cannot be
scored without it.
\item Repeat and machine-readable-output flags on the benchmark harness, which
currently reports only to the terminal and an exit code.
\item A bag-based offline scorer for E3--E5, in preference to a live scoring
node, so that runs can be re-analysed without re-execution.
\end{enumerate}

\subsection{Threats to validity}

\emph{Simulation only.} All results are in Stage 2D with simulator odometry as
ground truth and AMCL disabled. Localization error, wheel slip, and real sensor
noise are absent. Conclusions about orchestration transfer more readily than
conclusions about tracking accuracy.

\emph{The simulator pivot is qualitative.} The move from Gazebo to Stage was
driven by observed infeasibility at 20 robots, not by a controlled throughput
sweep. We report real-time factor as a function of fleet size and navigation
load for the Stage configuration only, and make no quantitative claim about the
crossover point. The same applies to bring-up reliability: staggered launch is
observed to work, but failure rate as a function of stagger interval and DDS
profile has not been measured over repeated cold boots.

\emph{Single host.} Twenty robots on one workstation is a DDS-bound ceiling;
timing metrics are contaminated by host load and should be read as relative,
not absolute.

\emph{Benchmark construction.} The 44 cases and the prompt were developed
together by the same authors, which risks the benchmark encoding the prompt's
idiosyncrasies. The category breakdown and the E1/E3 comparison partially
expose this, but we do not claim the benchmark is unbiased with respect to our
own system. An independent case set would be a genuine improvement.

\emph{Model non-stationarity.} Hosted models change under fixed names. We pin
local weights and report quantization for every reported run.

% ===========================================================================
\section{Results}

\subsection{E1: Planning accuracy}
\label{sec:e1-results}

Table~\ref{tab:e1-summary} covers all five models swept (Sec.~\ref{sec:e1}):
\textbf{qwen3.5:4b} (small), \textbf{gemma3:12b} (mid-size, different family),
\textbf{qwen2.5:14b} (4-bit, the intended operating point),
\textbf{qwen3.5-9b} (8-bit), and \textbf{mistral-small3.2:24b} (4-bit, largest
that fits). All runs use $N{=}5$ repeats per case (44 cases $\times$ 5 $=$ 220
calls per model), temperature 0.1, \code{--llm-max-tokens 2048},
\code{--llm-num-ctx 32768}.

\begin{table}[t]
\centering\footnotesize
\caption{E1 overall pass rates and latency (44 cases, $N{=}5$, 220 calls/model).}
\label{tab:e1-summary}
\begin{tabular}{lccccc}
\toprule
Model & pass@1 & pass@5 & med.\ (s) & p90 (s) & t/o \\
\midrule
qwen3.5:4b          & 68.2\% & 88.6\% &  4.8 &  8.6 & 0 \\
gemma3:12b           & 68.2\% & 70.5\% &  3.7 &  7.4 & 0 \\
qwen2.5:14b          & 72.7\% & 79.5\% &  9.3 & 20.1 & 0 \\
\textbf{qwen3.5-9b}  & \textbf{77.3\%} & \textbf{90.9\%} &  8.2 & 21.0 & 0 \\
mistral-sm3.2:24b     & 75.0\% & 77.3\% & 29.9 & 52.1 & 0 \\
\bottomrule
\end{tabular}
\end{table}

Raising the token budget from 768 to 2048 (Sec.~\ref{sec:e1}) removed timeouts
entirely across all five models. The best model on both metrics,
qwen3.5-9b, is not the largest: parameter count does not predict accuracy
here. Two further patterns stand out. First, the \emph{smallest} model,
qwen3.5:4b, is close behind on pass@1 and second-best on pass@5
(88.6\%)---the gap between its pass@1 and pass@5 is the largest in the table,
meaning most of its failures do not repeat under resampling and are closer to
sampling noise than to a structural gap. Second, gemma3:12b and
mistral-small3.2:24b show almost no such recovery (68.2\%$\to$70.5\% and
75.0\%$\to$77.3\%): their failures are largely the \emph{same} failure on
every retry, i.e.\ systematic rather than stochastic. Latency also does not
track accuracy: mistral-small3.2:24b's weights (15\,GB, 4-bit) leave under
1\,GB of headroom on the 16\,GB card, forcing partial CPU offload and a
median latency 6$\times$ qwen3.5-9b's---without a corresponding accuracy
gain.

The five models fail in qualitatively different ways
(Table~\ref{tab:e1-failuremodes}). \emph{Wrong formation fields}---almost
entirely a single validator complaint, \code{follower\_ns(N) !=
offsets\_x(M) or offsets\_y(M)}, i.e.\ the model emits an offsets array
shorter than its own follower list---is the single largest failure mode for
three of five models, and for mistral-small3.2:24b it is byte-identical
across 25 of its 57 failures. That repetition, combined with the model's
near-zero pass@1$\to$pass@5 recovery, is the signature of a deterministic
generation bug in how this model fills the formation-node schema, not of
noisy sampling. \emph{Wrong structure} (choosing the wrong top-level node
type, e.g.\ \code{mapf} instead of \code{sequence}) dominates for
qwen2.5:14b (28/60 failures) and gemma3:12b (25/68); qwen3.5:4b has no
single dominant mode, consistent with its noise-like failure profile above.
\emph{Wrong escalation reason} is essentially constant across four of the
five models ($\approx$10 failures each, qwen2.5:14b is the outlier at 5),
suggesting that judging when to ask for help rather than act is uniformly
hard and does not improve much with scale or family within this range---an
argument for a deterministic escalation check outside the model rather than
prompting harder.

\begin{table}[t]
\centering\footnotesize
\setlength{\tabcolsep}{2pt}
\caption{E1 failure-mode counts (failed calls only, out of 220 each).}
\label{tab:e1-failuremodes}
\begin{tabular}{lrrrrr}
\toprule
Failure mode & \makecell{qwen3.5\\4b} & \makecell{gemma3\\12b} & \makecell{qwen2.5\\14b} & \makecell{qwen3.5\\9b} & \makecell{mistral\\sm3.2 24b} \\
\midrule
Unparseable JSON    & 16 &  1 &  0 &  6 & 11 \\
Wrong structure      & 17 & 25 & 28 & 11 &  0 \\
Wrong robot-id set    &  0 &  0 &  0 &  0 &  1 \\
Wrong goals            & 11 &  0 & 10 &  8 &  5 \\
Wrong formation fields  & 22 & 24 &  9 &  2 & 25 \\
Wrong escalation reason  & 10 & 10 &  5 & 10 & 10 \\
Other validator          &  0 &  8 &  8 &  1 &  5 \\
\midrule
Total failed / 220        & 76 & 68 & 60 & 38 & 57 \\
\bottomrule
\end{tabular}
\end{table}

By category, all five models are reliable (100\% pass@1) on \code{idle} and
\code{parallel}, and near-reliable on \code{mapf\_basic} and \code{sequence}.
\code{formation\_create} shows the widest spread across models (100\% for
qwen3.5:4b down to 20\% for gemma3:12b), consistent with the formation-schema
failure mode above being model-specific rather than a shared difficulty. One
category is weak for \emph{every} model regardless of size or family:
\code{mapf\_spread} sits at exactly 50\% pass@1 for all five---worth flagging
as a possible benchmark-design artifact rather than a genuine model weakness,
since the category has only two cases and a single flipped call moves the
rate by 50 points; we do not read anything into it beyond that caveat.
\code{escalate} ranges 25--75\% pass@1 with no relationship to model size
(qwen2.5:14b, the largest of the mid-size models, is the \emph{worst} on this
category), reinforcing the failure-mode-table observation that escalation
judgment does not scale predictably.

\emph{Reading with appropriate caution.} $N{=}5$ per case, and as few as 2
cases in some categories, is enough to separate the reliable categories from
the weak ones but not enough to bound per-category confidence intervals
tightly; the category-level percentages above should be read as pointing at
where the failures cluster, not as precise estimates. We report them because
E2--E6 (Secs.~\ref{sec:e2}--\ref{sec:e6}) are designed to explain \emph{why}
those specific categories fail, not merely to confirm that they do.

\pending{E2 --- ablation: effect of constrained decoding, tool calling, pose
context, context budget.}

\pending{E3 --- end-to-end mission completion by configuration, with LLM-call
cost and guard/repair firing counts.}

\pending{E3b --- channel-1 decision confusion matrix and false-replan rate.}

\pending{E3c --- recovery and escalation under mid-mission door closure.}

\pending{E4 --- PBS vs.\ LNS2 across fleet size and clutter.}

\pending{E5 --- formation tracking error by shape and leader speed.}

% ===========================================================================
\section{Conclusion}

\pending{Full claim pending E2--E6. The intended claim remains: a small typed
plan grammar plus deterministic guards lets a modestly-sized locally-hosted
model supervise a 20-robot fleet, and the gap between plan-level accuracy and
mission-level success is where the engineering value sits---E3 is the
experiment that can actually support that sentence.}

What E1 alone supports is narrower but not nothing. First, bare planning
accuracy without any of the system's guards or repair loops sits at
68--77\% pass@1 across five locally-hosted 3.4--24B models under a realistic
$\approx$9{,}700-token prompt---a fleet-supervision system built on a single
unguarded plan call would fail roughly one command in three to one in four
\emph{regardless of which model is deployed}, which is the quantitative case
for the defence-in-depth design of Sec.~\ref{sec:guards} rather than an
assumption behind it. Second, the five models' failures are not
interchangeable, and not even uniformly noisy: the smallest model's errors
look like sampling noise (pass@5 recovers most of the gap to a near-perfect
pass@1), while two mid/large models repeat the \emph{identical} validator
complaint---a malformed \code{offsets\_x}/\code{offsets\_y} array in
formation nodes---on retry after retry. That argues for model-specific
schema-repair logic rather than a single fixed remediation budget applied
uniformly. Third, escalation judgment (recognizing missing or ambiguous
state and asking for help instead of acting) is weak and essentially flat
across all five models independent of size or family, which is the
strongest single piece of evidence in this data for keeping that decision
outside the model rather than trying to prompt it better. Fourth, the
categories that are hardest across the board---formation staging order,
escalation judgment---are exactly the categories the plan grammar
(Sec.~\ref{sec:planlang}) and the guard layer were built to catch
downstream, which is a consistency check on the architecture rather than a
result about it. Whether the guards actually close that gap in a running
fleet, and whether it holds under tool calling instead of constrained
decoding, is what E2 and E3 are for.

\section*{Code availability}

\url{https://github.com/innopolis-robotics-society/llm_swarm}

% ===========================================================================
\begin{thebibliography}{99}\small
\bibitem{ourlevel2} E. Shlomov, A. Malakhova, S. Belyh, M. Shariev,
``Hierarchical Multi-Agent Control by LLM: Level-2 Swarm Coordination via MAPF,
Formation Control, and Behavior Trees,'' technical report, 2026.
\bibitem{stern} R. Stern et al., ``Multi-agent pathfinding: Definitions,
variants, and benchmarks,'' SoCS, 2019.
\bibitem{cbs} G. Sharon et al., ``Conflict-based search for optimal multi-agent
pathfinding,'' \emph{Artificial Intelligence}, 2015.
\bibitem{pbs} H. Ma et al., ``Searching with consistent prioritization for
multi-agent path finding,'' AAAI, 2019.
\bibitem{lns2} J. Li et al., ``MAPF-LNS2: Fast repairing for multi-agent path
finding via large neighborhood search,'' AAAI, 2022.
\bibitem{formsurvey1} C. Jiang, Z. Chen, and Y. Guo, ``Multi-robot formation
control: a comparison between model-based and learning-based methods,''
\emph{J. Control Decis.}, vol. 7, no. 1, pp. 90--108, 2020.
\bibitem{formsurvey2} E. Debie, K. Kasmarik, and M. Garratt, ``Swarm robotics:
A survey from a multi-tasking perspective,'' \emph{ACM Comput. Surv.}, vol. 56,
no. 2, 2023.
\bibitem{apf} H. Wang, M. Yue, X. Sun, and X. Zhao, ``Cooperative formation
control strategy for multi-robot system based on APF algorithm and sliding mode
estimator,'' \emph{Robot. Auton. Syst.}, vol. 193, p. 105075, 2025.
\bibitem{nav2} S. Macenski et al., ``The Marathon 2: A navigation system,''
IROS, 2020.
\bibitem{layered} D. V. Lu et al., ``Layered costmaps for context-sensitive
navigation,'' IROS, 2014.
\bibitem{btbook} M. Colledanchise and P. Ögren, \emph{Behavior Trees in
Robotics and AI}, CRC Press, 2018.
\bibitem{btcpp} D. Faconti, ``BehaviorTree.CPP.'' \url{https://www.behaviortree.dev}
\bibitem{saycan} M. Ahn et al., ``Do as I can, not as I say: Grounding language
in robotic affordances,'' CoRL, 2022.
\bibitem{codeaspolicies} J. Liang et al., ``Code as policies: Language model
programs for embodied control,'' ICRA, 2023.
\bibitem{progprompt} I. Singh et al., ``ProgPrompt: Generating situated robot
task plans using large language models,'' ICRA, 2023.
\bibitem{innermonologue} W. Huang et al., ``Inner monologue: Embodied reasoning
through planning with language models,'' CoRL, 2022.
\bibitem{dds1} Y. Maruyama, S. Kato, and T. Azumi, ``Exploring the performance
of ROS2,'' EMSOFT, 2016.
\bibitem{dds2} T. Kronauer, J. Pohlmann, M. Matthe, T. Smejkal, and
G. Fettweis, ``Latency analysis of ROS2 multi-node systems,'' 2023.
\bibitem{ros2} S. Macenski, T. Foote, B. Gerkey, C. Lalancette, and W. Woodall,
``Robot Operating System 2: Design, architecture, and uses in the wild,''
\emph{Sci. Robot.}, vol. 7, no. 66, 2022.
\bibitem{stage} R. Vaughan, ``Massively multi-robot simulation in Stage,''
\emph{Swarm Intell.}, vol. 2, pp. 189--208, 2008.
\bibitem{gazebo} N. Koenig and A. Howard, ``Design and use paradigms for
Gazebo, an open-source multi-robot simulator,'' IROS, 2004.
\bibitem{btsurvey} M. Iovino, E. Scukins, J. Styrud, P. Ögren, and C. Smith,
``A survey of behavior trees in robotics and AI,'' \emph{Robot. Auton. Syst.},
vol. 154, 2022.
\bibitem{toolformer} T. Schick et al., ``Toolformer: Language models can teach
themselves to use tools,'' NeurIPS, 2023.
\bibitem{llmrobsurvey} F. Zeng et al., ``Large language models for robotics: A
survey,'' 2023.
\bibitem{rt2} A. Brohan et al., ``RT-2: Vision-language-action models transfer
web knowledge to robotic control,'' arXiv:2307.15818, 2023.
\bibitem{secondteam} D. Vizitei, A. Tuzov, D. Ryabov, ``Hierarchical
Multi-Agent Control by LLM: Simulation Environment, Per-Robot Navigation, and
LLM Orchestration,'' technical report, 2026.
\end{thebibliography}

\end{document}
